\section{Policy implications}\label{sec:policy}

\paragraph{Pre-trade controls.} Rule 15c3-5 requires pre-trade controls that prevent orders exceeding credit or capital thresholds \citep{sec2010}. Static single-order caps and daily value-at-risk limits do not see transient dynamics. Lemma~\ref{thm:guard} identifies the quantity that a solvency-motivated control should bound: inventory times the shock envelope times the summed absolute impulse response. Peak-based or eigenvalue-based limits are either unsafe or blind in exactly the regimes where they matter. Our evidence locates where such controls matter. In the deepest crypto book the amplifying regime is rare and the statistic carries no forecasting value; in thin books it is more frequent, and six-hour estimates of the summed response forecast the next hour's adverse excursion beyond volatility forecasts and price impact (Section~\ref{sec:thin}). Our backtests also set a limit on what such a control can do with public data. A six-hour estimate of the summed response adds nothing to a volatility limit as a sizing input (Table~\ref{tab:backtest}); an operator certified from three days of data improves a volatility limit modestly, by 2.4\% in tail loss per unit of position (Table~\ref{tab:certguard}); and limits built on the estimate alone do worse than volatility in both. The guard's role is therefore a ceiling on worst-case loss under a certified operator and a correction to conventional volatility limits, not a replacement for them, and it is most relevant for books with under about \$0.4M of depth within 1\% of the mid.

\paragraph{Pricing reversion.} Proposition~\ref{prop:equil} shows that dealers under-supply reversion because its benefit, lower tail exposure for every leveraged position, is not paid to them. Maker rebates are a natural instrument: a rebate that rises with the book's summed response, $(1-\theta/N)\Lambda S^2/\tilde\kappa$ per unit of reversion, prices the externality with an observable. Venues already tier rebates by volume; tiering them by contribution to reversion in thin books is a design question our model makes precise.

\paragraph{What can be certified, what cannot, and when to abstain.} A supervisor or venue can certify the summed response, and hence a safe admission level, from controlled experiments at a cost that does not depend on how many traders there are or how similar their memories are (Theorem~\ref{thm:certcost}). It cannot certify any finite risk bound from a finite window without a bound on how long participants remember (Theorem~\ref{thm:window}(iii)), so a declared memory bound is a policy input, not a technicality, and a certificate must be renewed when conditions change (test T09). When compatible risks straddle the budget, the right response is to abstain: throttle admission to the certified level, and when even zero admission is uncertified because of existing inventory, flag the state and reduce the position (Theorem~\ref{thm:robustguard}(iv)). Disclosure is worth paying for where memory is long and public certificates are weak, and is worth little where a short window already pins the risk down (Theorem~\ref{thm:disclosure}).

\paragraph{Market design.} Frequent batch auctions remove latency races within a batch \citep{budish2015}. Transient amplification, when it exists, operates over many seconds and spans many batches, so batching and a state-carrying guard address different failures and can be complements.

\paragraph{Supervision without identification.} The companion paper shows that identifying algorithms from public data is expensive. The guard needs only impulse coefficients and operator confidence sets, which are invariant to how flow is split across cohorts. Supervisors can therefore ask venues and brokers to report summed-gain statistics and their confidence radii instead of attempting to infer black-box learners. Our results argue for reporting impact, mean reversion and refill directly, because those are the primitives that decide whether a book can amplify. They also show that feedback and amplification are different objects: measured feedback is largest in the deepest books, where amplification is rarest (Table~\ref{tab:x3}).

\section{Conclusion}\label{sec:conclusion}

The theory answers four questions. A book amplifies when depletion outweighs impact-weighted reversion; with momentum the one-period margin is sufficient but not exact, and momentum alone cannot make a stable book overshoot. Cumulative risk is long-run impact divided by reversion net of momentum, and the guard that bounds it is the largest safe pre-trade rule, computable and identified from public data. In equilibrium reversion is under-supplied, and the corrective price is a function of the guard's own statistic. And the information that certifies a safe decision is much cheaper than the information that identifies the traders: the summed response is certified at the parametric rate up to a logarithm that cannot be removed, finite windows certify it only under a memory bound, and a guard built on the certificate gives up no more trading capacity than any comparably safe rule, up to a logarithm.

The evidence is mixed in an informative way. The deepest crypto book sits on the stable side of the condition, and flow depletes the side it hits without producing a peak overshoot. In thin Binance books, six-hour estimates of cumulative transient risk forecast the size of the next hour's adverse excursion out of sample in every pre-registered forecasting test, on new symbols and on untouched days, beyond price impact, a full volatility-forecasting baseline and VPIN; the effect is about a fifth as large in thin US stocks, negligible on Bybit and absent on Coinbase. Tests of a leverage or liquidation mechanism fail, and a follow-up that points to slow reversion over 15 to 60 minutes is close to a check of how the statistic is built. As a risk control the certified guard improves on a static limit but not, at matched exposure, on a volatility cap; its registered gain over a volatility cap in 18-hour positions comes mostly from admitting less exposure. Thin books are risky because their prices revert too slowly relative to the impact and momentum of the flow that moves them; why they revert slowly is the open empirical question.

Three extensions would settle what remains open. Message-level data from a regulated market, such as LOBSTER or TAQ for US equities, would extend the US test to depth and to large stocks, show where the effect ends, and allow depletion to be measured without the recentering of band depth. Data on liquidity-provider identities would test the equilibrium's prediction that fewer dealers mean more transient risk. And longer certification windows, or certification with order-level data, would show whether the modest gain of a certified guard over volatility limits can be made large.

\paragraph{Data and code availability.} All data are public: Binance archives (\url{https://data.binance.vision}), verified against published checksums, and Bybit's order-book history (\url{https://www.bybit.com/derivatives/en/history-data}) and trade archive (\url{https://public.bybit.com}); US consolidated trades are from the Alpaca market-data API (\url{https://alpaca.markets}, free account). Pre-registered specifications, their SHA-256 hashes and test-set unlock counters are in the replication package. Each number is produced by a notebook cell or script that writes a result file (Appendix~\ref{app:repro}).

\paragraph{Use of AI tools.} The authors used Claude (Anthropic), an AI assistant, to write and debug analysis code, run checks of derivations and numerical results, and draft and edit text. Every number in the paper is produced by code the authors executed and archived in the replication package. The authors reviewed all content and take full responsibility for it.

\paragraph{Disclosures.} The authors declare no relevant financial interests.
